\section{Methodology and Test Framework}
\label{sec:methodology}

The benchmark compares 16 frequency estimators on a common dynamic test set in
a local, non-time-synchronized measurement setting, using waveform generation
at $1$\,MHz and estimator outputs at $10$\,kHz. The scenario set extends the
same protection-oriented stress modes used in standard dynamic
testing--magnitude steps, ramps, modulation, and phase jumps--with
IBR-specific harmonic, out-of-band, ringdown, and multi-event stress. This
expansion does not change the measurement setting: all methods are evaluated as
local relay-class frequency trackers from the same decimated $10$\,kHz data
stream.

The estimator pool spans six practical groups: zero-crossing and energy
baselines (ZCD, TKEO), windowed spectral methods (IpDFT, TFT), PLL/SOGI-based
trackers (PLL, SOGI-PLL, SOGI-FLL, Type-3 SOGI-PLL), Kalman-filter methods
(LKF2, EKF, UKF, RA-EKF), parametric/subspace methods (Prony, ESPRIT), and
reduced-order or adaptive/data-driven designs (RLS, Koopman/RK-DPMU). This mix
is intended to expose trade-offs relevant to protection and monitoring:
accuracy, transient risk, settling behavior, and computational burden. Within
the Kalman-filter group, EKF and UKF use a single-tone state over phase,
frequency, and amplitude, whereas RA-EKF adds an explicit RoCoF state,
innovation-scaled process noise, and a soft phase reset under large
innovations. RoCoF augmentation in EKF-type filters has prior literature
antecedents~\cite{Ferrero2020EKF,Chamorro2021}; the benchmark treats RA-EKF as
one Kalman-family variant, not as a new estimator family.

The study contains 34 benchmark scenarios. Rather than read the benchmark as a
flat scenario list, we group the scenarios by disturbance family.
Table~\ref{tab:benchmark_scope} summarizes the resulting ten-family structure,
and Fig.~\ref{fig:scenario_overview} shows representative ground-truth signals
for six of those families.

\begin{table}[t]
  \centering
  \caption{Benchmark scope. Counts refer to benchmark scenarios. OOB denotes
    out-of-band interference. The phase-jump family includes $20^\circ$ and
    $60^\circ$ re-lock tests.}
  \label{tab:benchmark_scope}
  \scriptsize
  \setlength{\tabcolsep}{3pt}
  \renewcommand{\arraystretch}{1.08}
  \begin{tabular}{L{1.7cm} c L{3.95cm}}
    \toprule
    Family & Count & Stress tested \\
    \midrule
    IEEE ramps & 9 & Sustained RoCoF from 0.25 to 20\,Hz/s \\
    Mag.\ steps & 7 & 1\% to 50\% magnitude changes \\
    Modulation & 3 & Nominal, AM, and FM tracking \\
    Phase jumps & 3 & $20^\circ$ and $60^\circ$ phase-jump re-lock tests \\
    OOB interference & 1 & Spectral rejection under off-nominal content \\
    Single-tone & 1 & Deterministic steady-state control case \\
    IBR harmonics & 3 & Small, medium, and large harmonic stress \\
    IBR ringdown & 5 & Power-imbalance ringdown with noise sweep \\
    IBR multi-event & 1 & Composite phase jump, ramp, harmonics, and ringdown \\
    Freq.\ step & 1 & Isolated frequency-step control case \\
    \bottomrule
  \end{tabular}
\end{table}

The family structure reflects distinct physical stress modes. Ramps probe
sustained RoCoF tracking. Magnitude steps and modulation probe
frequency-amplitude decoupling. Phase jumps probe re-lock behavior and
transient overshoot. Out-of-band interference and IBR harmonics probe spectral
leakage. Ringdown and multi-event cases combine these effects under changing
noise and transient recovery, which makes them the strongest stress tests in
the benchmark. Two records drive the sharpest separations: a NERC-type
$+60^\circ$ phase-jump case at $t=0.7$\,s with 4\% 5th, 2\% 7th, and 0.5\%
32.5-Hz interharmonic content, and an IBR multi-event case with a 0.5-pu
sag/swell, a $+45^\circ$ phase jump, a $-2$\,Hz/s RoCoF ramp clamped at
58\,Hz, decaying DC, AM/PM oscillations, a full IBR harmonic spectrum, and
compound noise. Figure~\ref{fig:scenario_overview} illustrates that
difference directly: some benchmark families are isolated compliance-style
events, whereas others stack abrupt discontinuities, ramps, and oscillatory
recovery inside the same record.

\begin{figure*}[t]
  \centering
  \includegraphics[width=0.98\textwidth,height=0.45\textheight,keepaspectratio]{Fig1_Scenarios_Final.pdf}
  \caption{Representative benchmark scenarios. Each row shows the voltage
    waveform and true frequency trajectory for six disturbance families:
    10\% magnitude step, frequency ramp, modulation, $60^\circ$ phase jump,
    IBR multi-event sequence, and IBR power-imbalance ringdown. Dashed lines
    mark disturbance onset. The figure illustrates why the benchmark is read by
    disturbance family: some cases isolate standard dynamic stress, whereas
    others combine discontinuities, ramps, harmonics, and oscillatory recovery
    in one record. Panel (E) shows the dominant frequency profile of the
    multi-event sequence; the full scenario also includes a 0.5-pu sag/swell,
    AM/PM modulation, a full IBR harmonic spectrum, and compound noise.}
  \label{fig:scenario_overview}
\end{figure*}

All estimators are evaluated as local relay-class frequency trackers from the
same decimated $10$\,kHz data stream; the study does not assess PMU
time-stamping or synchrophasor compliance. Each estimator--scenario pair is
tuned via scenario-wise parameter search across 100 trials and then evaluated
over 60 Monte Carlo runs. The reported values therefore approximate
best-in-range performance within the tested search space, not untuned default
settings. The study contains 544
estimator--scenario summaries derived from 32,640 Monte Carlo records. Let
$e[k]=\hat{f}[k]-f_{\mathrm{true}}[k]$. We report
$\mathrm{RMSE}=\sqrt{N^{-1}\sum_k e[k]^2}$, maximum peak error
$\max_k |e[k]|$, trip-risk duration
$T_{\mathrm{trip}}=\Delta t\sum_k \mathbf{1}\{|e[k]|>0.5\,\mathrm{Hz}\}$, settling time,
and CPU time per sample. The 0.5-Hz threshold is used as a
protection-oriented risk proxy rather than as a universal trip criterion. CPU
time is interpreted as a relative computational-cost indicator measured in one
software environment, not as hardware latency. For any metric $m$, the family
summary is $\bar{m}_F=|F|^{-1}\sum_{s\in F}\bar{m}_s$, where $\bar{m}_s$ is
the scenario mean over the 60 realizations. This prevents dense families from
dominating the comparison. Unless stated otherwise, frequency errors are
reported in mHz in the tables and discussion.
